\section{需求分析}

\subsection{用户角色}

本平台面向四类核心用户角色：

\begin{table}[H]
  \centering
  \caption{用户角色分析}
  \label{tab:user-roles}
  \begin{tabularx}{\textwidth}{p{2.5cm}|p{4.5cm}|p{4.5cm}|p{3.5cm}}
    \toprule
    \textbf{角色} & \textbf{特征} & \textbf{核心需求} & \textbf{平台价值} \\
    \midrule
    钓鱼爱好者（新手） & 对钓鱼感兴趣但经验不足，依赖工具决策 & 可视化地图找点、智能推荐、系统化教程、良好社区氛围 & 降低入门门槛，缩短从“想去钓”到“会钓”的距离 \\
    \hline
    钓鱼爱好者（老手） & 具有稳定出钓频率，关注钓点质量和社群认可 & 钓点质量判断、鱼种/鱼情/钓法信息、装备讨论、战绩展示 & 提供结构化数据沉淀和社群认同感 \\
    \hline
    渔业经营主体 & 钓场、渔庄等，需要展示和获客 & 商户展示、预约转化、赛事组织、用户评价管理 & 精准触达垂钓消费人群 \\
    \hline
    地方监管与服务主体 & 渔政、文旅、农业农村部门 & 禁渔区宣传、文明垂钓引导、违规提醒、产业数据监测 & 提供数字化治理抓手 \\
    \bottomrule
  \end{tabularx}
\end{table}

当前 MVP 阶段优先服务前两类用户（钓鱼爱好者），渔业经营主体和监管服务主体作为后续扩展方向。

\subsection{功能需求}

围绕五个核心用户问题，平台需提供以下功能模块：

\begin{enumerate}
  \item \textbf{地图与 POI 发现}：展示附近钓点、支持关键字搜索与类型筛选（野钓/钓场/海钓/免费/夜钓/路亚）、展示天气摘要、禁钓区与风险图层开关。当前 MVP 阶段基于高德 JS API 实现地图展示，POI 数据以 JSON 集合提供，后续可接入高德 POI 搜索和自建垂钓 POI。
  \item \textbf{POI 详情与推荐}：展示钓点名称、类型、距离、鱼种、收费、设施、规则摘要；基于鱼情分、合规分、距离分、天气适配分、设施便利分、社交热度分计算推荐指数；对禁钓、限钓、极端天气等风险优先提示。
  \item \textbf{出钓记录}：支持鱼获记录、空军记录、探点记录三种类型；自动回填天气、时间、位置；结构化采集鱼种、数量、重量、钓法、饵料、装备等字段；生成可分享的战绩卡片。
  \item \textbf{社区与发布}：支持鱼获流、同城流和话题流；帖子关联出钓记录；内容权限（公开/私密）与位置权限（精确/模糊/隐藏）分离。
  \item \textbf{AI 推荐解释}：后端规则模型完成排序，DeepSeek 将结构化推荐结果解释为自然语言建议，不直接参与排序决策。
  \item \textbf{教程学习}：按新手入门、钓法教程、鱼种攻略、场景攻略等分类组织；教程关联附近练习钓点。
  \item \textbf{个人主页与统计}：展示总出钓、总鱼获、常钓鱼种、高频钓点、偏好时段等统计；当前 MVP 阶段已提供前端展示和后端 profile-summary/report 汇总接口，后续可扩展更完整的月报、年报和会员报告。
\end{enumerate}

\subsection{非功能需求}

\begin{table}[H]
  \centering
  \caption{非功能需求}
  \label{tab:nfr}
  \begin{tabularx}{\textwidth}{p{3cm}|p{6cm}|p{6cm}}
    \toprule
    \textbf{类别} & \textbf{需求描述} & \textbf{MVP 阶段策略} \\
    \midrule
    性能 & 页面首次加载 < 3 秒，地图 POI 渲染流畅 & 前端地图组件动态加载，主 JS 包控制体积 \\
    \hline
    可用性 & 未授权定位时引导授权并提供手动城市选择；API 异常时展示缓存数据和错误提示 & 前端统一错误处理，避免静默降级为假数据 \\
    \hline
    安全性 & API Key 不写入前端代码，用户数据不得通过任意 user\_id 访问 & 第三方 Key 由后端环境变量管理，后续增加基础鉴权 \\
    \hline
    可扩展性 & 前后端分离，接口字段稳定，新增数据集合同步创建 .schema.json & 现有三层架构（路由 → Service → JSON Store）支持逐步演进 \\
    \hline
    兼容性 & 支持微信小程序与 iOS/Android 移动端 & 当前 MVP 先以 Web 移动端完成演示，Uni-app 跨端作为后续方向 \\
    \hline
    数据持久性 & 新增记录和帖子在服务重启后不丢失 & JSON 文件原子读写（临时文件 + 原子替换） \\
    \bottomrule
  \end{tabularx}
\end{table}

\subsection{竞品分析}

当前垂钓相关产品呈现“多而分散”的格局，用户完成一次出钓决策通常需要切换多个平台：

\begin{table}[H]
  \centering
  \caption{竞品格局分析}
  \label{tab:competitors}
  \begin{tabularx}{\textwidth}{p{2.5cm}|p{4cm}|p{4cm}|p{4.5cm}}
    \toprule
    \textbf{类型} & \textbf{代表产品} & \textbf{优势} & \textbf{不足} \\
    \midrule
    传统地图平台 & 高德地图、百度地图 & 底图能力强、POI 检索成熟、导航体验好 & 缺乏垂钓语义层，无法回答“能不能钓、有没有鱼” \\
    \hline
    垂直钓鱼社区 & 钓鱼人、钓鱼之家 & 社区氛围强、内容深度好、用户信任度高 & 地图表达弱、政策整合不足、数据标准化程度低 \\
    \hline
    泛内容平台 & 抖音、小红书、B站 & 流量巨大、内容分发能力强、适合破圈传播 & 地点可信度弱、信息结构化程度低、难以沉淀为空间数据资产 \\
    \hline
    工具类平台 & 墨迹天气、两步路、中国海洋预报 & 单点问题解决能力强、用户理解成本低 & 无法覆盖完整出钓链路（找点→判断→记录→分享→复盘） \\
    \bottomrule
  \end{tabularx}
\end{table}

「鱼你有图」的定位不是再造一个地图或论坛，而是在成熟地图底图之上叠加垂钓专题 POI、禁钓规则、鱼情数据和钓友内容，形成差异化的垂钓 GIS 应用层；同时把泛内容平台作为内容引流阵地，承接用户从“看到鱼获”到“查询钓点、判断规则、规划出行、记录鱼获”的后续决策链路。
